\documentclass[11pt]{article}

\usepackage[margin=1in]{geometry}
\usepackage{amsmath,amssymb,amsthm}
\usepackage{graphicx}
\usepackage[hidelinks]{hyperref}

\newtheorem{theorem}{Theorem}
\newtheorem{lemma}{Lemma}
\newtheorem{remark}{Remark}

\newcommand{\N}{\mathbb{N}}
\newcommand{\R}{\mathbb{R}}
\newcommand{\co}{\operatorname{co}}
\newcommand{\supp}{\operatorname{supp}}
\newcommand{\norm}[1]{\left\lVert #1 \right\rVert}
\newcommand{\eps}{\varepsilon}

\title{Initial Sign Vectors in \texorpdfstring{$c_0$}{c0} Refute a Schur-Generation Question}
\author{Run \texttt{fa\_banach\_001}, agent \texttt{agent\_lane\_01}}
\date{June 16, 2026}

\begin{document}
\maketitle

\section*{Source Question}

The source is Ondrej Kurka,
\emph{Tsirelson-like spaces and complexity of classes of Banach spaces},
arXiv:1612.09334.  Remark 5.9(iii), on page 21, asks:

\begin{quote}
Does a Banach space \(X\) has necessarily the Schur property if it has a
subset \(W\) such that \(\overline{\co} W=B_X\) and every weakly convergent
sequence of elements of \(W\) is convergent in the norm?
\end{quote}

\begin{center}
\includegraphics[width=0.98\linewidth]{figures/open_problem_crop.png}
\end{center}

\section*{Answer}

The answer is negative.  The counterexample is \(X=c_0\), with \(W\) the
union of all finite initial sign cubes.

\section*{Definitions and Normalizations}

Let \(c_0\) be the real Banach space of scalar sequences converging to
zero, with the supremum norm.  The Schur property means that every weakly
convergent sequence is norm convergent.

For \(n\geq 1\), define
\[
  W_n=\left\{
  (\eps_1,\ldots,\eps_n,0,0,\ldots):
  \eps_i\in\{-1,1\}\ \text{for }1\leq i\leq n
  \right\}\subset c_0
\]
and put
\[
  W=\bigcup_{n=1}^{\infty} W_n.
\]
Every element of \(W\) belongs to the unit sphere of \(c_0\).

\section*{Proof Intuition}

The set \(W_n\) is the vertex set of the \(n\)-dimensional cube
\([-1,1]^n\), embedded in the first \(n\) coordinates of \(c_0\).  Thus
the convex hull of \(W\) contains every finitely supported point of the
unit ball, and these points are dense in \(B_{c_0}\).

The weak-sequential condition holds for a different reason.  A weakly
convergent sequence in \(c_0\) is coordinatewise convergent.  If the sign
vectors have unbounded lengths, then after passing to a subsequence every
fixed coordinate is eventually a sign, and coordinate convergence forces
that coordinate of the weak limit to have modulus one.  This would happen
for every coordinate, which is impossible for an element of \(c_0\).  So
the lengths are bounded, and the sequence lives in a finite set.

\section*{Counterexample}

\begin{theorem}
There is a Banach space \(X\) and a subset \(W\subset B_X\) such that
\[
  \overline{\co} W=B_X,
\]
and every weakly convergent sequence of elements of \(W\) is norm
convergent, but \(X\) does not have the Schur property.
\end{theorem}

\begin{proof}
Take \(X=c_0\) and \(W=\bigcup_{n\geq 1} W_n\) as above.

First we show that \(\overline{\co} W=B_{c_0}\).  For each fixed \(n\),
the convex hull of \(W_n\) is exactly the finite cube
\[
  \{x\in c_0:\supp x\subseteq\{1,\ldots,n\},\ |x_i|\leq 1
  \text{ for }1\leq i\leq n\}.
\]
Indeed, this is the usual description of the cube as the convex hull of
its \(2^n\) sign vertices.  Hence every finitely supported \(x\in B_{c_0}\)
belongs to \(\co W\), after choosing \(n\) with \(\supp x\subseteq
\{1,\ldots,n\}\).  Since the finitely supported points of \(B_{c_0}\) are
dense in \(B_{c_0}\) by coordinate truncation, \(\overline{\co} W=B_{c_0}\).

Now let \((w_j)_j\) be a weakly convergent sequence in \(W\), say
\[
  w_j \xrightarrow{\ \mathrm{weak}\ } x\in c_0.
\]
Let \(n_j\) be the unique integer such that \(w_j\in W_{n_j}\).  We claim
that the sequence \((n_j)_j\) is bounded.  Suppose not.  Passing to a
subsequence, still denoted \((w_j)_j\), we may assume that \(n_j\to\infty\).
For each fixed coordinate \(k\), the coordinate functional
\(\delta_k\in c_0^*=\ell_1\) gives
\[
  \delta_k(w_j)\longrightarrow \delta_k(x)=x_k.
\]
But for all sufficiently large \(j\), since \(k\leq n_j\), we have
\(\delta_k(w_j)\in\{-1,1\}\).  A convergent scalar sequence whose tail is
contained in \(\{-1,1\}\) is eventually constant.  Therefore
\[
  |x_k|=1
\]
for every \(k\in\N\).  This contradicts \(x\in c_0\).  Thus the lengths
\((n_j)_j\) are bounded.

Choose \(N\) such that \(n_j\leq N\) for all \(j\).  Then
\[
  \{w_j:j\in\N\}\subseteq \bigcup_{n=1}^N W_n,
\]
which is a finite set.  Since \((w_j)_j\) is weakly convergent, every
subsequence has the same weak limit.  A finite-valued convergent sequence
in the Hausdorff weak topology must eventually be constant: otherwise two
distinct values would occur on infinite subsequences and both would have
to equal the weak limit.  Hence \((w_j)_j\) is eventually constant, and in
particular norm convergent.

Finally, \(c_0\) does not have the Schur property.  Let \((e_n)_n\) be the
canonical unit vector basis.  For every \(a=(a_n)_n\in\ell_1=c_0^*\),
\[
  a(e_n)=a_n\longrightarrow 0,
\]
so \(e_n\to 0\) weakly in \(c_0\).  However, \(\norm{e_n}_\infty=1\) for
all \(n\), so \((e_n)_n\) is not norm convergent to \(0\).  Therefore
\(c_0\) fails the Schur property, while the set \(W\) satisfies the two
conditions in the source question.
\end{proof}

\section*{Verification Notes}

The proof is self-contained modulo the standard identification
\(c_0^*=\ell_1\).  The construction also survives harmless variations of
the wording in the source question: since \(\overline{\co} W=B_X\), any
such \(W\) must in particular lie inside \(B_X\), and the chosen \(W\)
already lies in the unit sphere.

The only sequential point is the bounded-length claim.  It uses
coordinatewise convergence, forced by weak convergence in \(c_0\), and
the fact that a vector whose every coordinate has modulus one cannot
belong to \(c_0\).

\section*{Novelty and Duplicate Checks}

Before promotion, I searched the run indexes for arXiv \texttt{1612.09334},
the title terms, \texttt{Tsirelson-like spaces}, \texttt{Schur property},
\texttt{weakly convergent sequence of elements of W}, and related phrases
from Remark 5.9(iii).  No existing packet for this exact question was
found; the nearby hits concerned different Schur, Tsirelson, or Borel
complexity questions.

I also ran bounded web searches for exact and near-exact phrases including
\begin{itemize}
\item \texttt{"every weakly convergent sequence of elements of W" "Schur property"},
\item \texttt{"Does a Banach space X" "Schur property" "subset W"},
\item \texttt{"Kurka" "subset W" "Schur property"},
\item \texttt{"c\_0" "every weakly convergent sequence" "norm convergent" "closed convex hull"}.
\end{itemize}
The searches did not surface a prior answer to this embedded question.
The packet is therefore marked as a candidate original counterexample
pending human review.

\section*{References}

\begin{thebibliography}{9}

\bibitem{Kurka}
O. Kurka,
\emph{Tsirelson-like spaces and complexity of classes of Banach spaces},
arXiv:1612.09334, 2016.

\end{thebibliography}

\section*{Human Review Recommendation}

This should be reviewed as an elementary counterexample to the displayed
question in Remark 5.9(iii).  The suggested checks are: verify the
transcription of the question, confirm that no additional hidden
restriction on \(W\) is intended, and inspect the weak-convergence
argument for the initial sign-vector set in \(c_0\).  Under the displayed
wording, the answer is decisively negative.

\end{document}
